\documentclass[11pt]{article}
\usepackage[margin=1in]{geometry}
\usepackage{amsmath,amssymb,mathtools}
\usepackage{enumitem}
\setlist[enumerate]{leftmargin=*,itemsep=0.25em,topsep=0.4em}

\title{Solving Linear Systems with QR and Cholesky Factorizations}
\author{}
\date{}

\begin{document}
\maketitle

\section{Solving a Linear System Using QR Decomposition}

Consider the system
\[
Ax=b.
\]
The QR method proceeds as follows.
\begin{enumerate}
  \item Compute the QR decomposition $A=QR$.
  \item The columns of $Q$ are orthonormal, so $Q^TQ=I$.
  \item The matrix $R$ is upper triangular.
  \item Substitute $A=QR$ into the system to obtain $QRx=b$.
  \item Multiply by $Q^T$: $Q^TQRx=Q^Tb$.
  \item Since $Q^TQ=I$, this becomes $Rx=Q^Tb$.
  \item Define $y=Q^Tb$.
  \item The remaining system is the upper-triangular system $Rx=y$.
  \item Solve it for $x$ by back substitution.
  \item If $A$ is rectangular with full column rank, the same procedure gives the least-squares solution minimizing $\lVert Ax-b\rVert_2$.
\end{enumerate}

\section{A Cholesky Example}

Solve
\[
\begin{bmatrix}4&2\\2&3\end{bmatrix}
\begin{bmatrix}x_1\\x_2\end{bmatrix}
=
\begin{bmatrix}6\\7\end{bmatrix}.
\]
Let $A=LL^T$, where
\[
L=\begin{bmatrix}l_{11}&0\\l_{21}&l_{22}\end{bmatrix}.
\]
Then
\[
LL^T=
\begin{bmatrix}
l_{11}^2 & l_{11}l_{21}\\
l_{11}l_{21} & l_{21}^2+l_{22}^2
\end{bmatrix}.
\]
Comparing entries with $A$ gives
\[
l_{11}=2,\qquad l_{21}=1,\qquad l_{22}=\sqrt{2},
\]
and hence
\[
L=\begin{bmatrix}2&0\\1&\sqrt{2}\end{bmatrix}.
\]
First solve $Ly=b$:
\[
2y_1=6,\qquad y_1+\sqrt{2}y_2=7,
\]
so
\[
y_1=3,\qquad y_2=2\sqrt{2}.
\]
Next solve $L^Tx=y$:
\[
\sqrt{2}x_2=2\sqrt{2},\qquad 2x_1+x_2=3.
\]
Thus
\[
x_2=2,\qquad x_1=\frac12,
\qquad
x=\begin{bmatrix}\frac12\\2\end{bmatrix}.
\]
Indeed,
\[
Ax=
\begin{bmatrix}4(1/2)+2(2)\\2(1/2)+3(2)\end{bmatrix}
=\begin{bmatrix}6\\7\end{bmatrix}=b.
\]

\section{The Cholesky Method in Algebraic Form}

Let $A\in\mathbb{R}^{n\times n}$ be symmetric positive definite.  To solve $Ax=b$:
\begin{enumerate}
  \item Compute the Cholesky factorization $A=LL^T$, where $L$ is lower triangular.
  \item For $j=1,\ldots,n$, compute the diagonal entries
  \[
  l_{jj}=\sqrt{a_{jj}-\sum_{k=1}^{j-1}l_{jk}^2}.
  \]
  \item For $i=j+1,\ldots,n$, compute the entries below the diagonal
  \[
  l_{ij}=\frac{a_{ij}-\sum_{k=1}^{j-1}l_{ik}l_{jk}}{l_{jj}}.
  \]
  \item Substitute the factorization into the system: $LL^Tx=b$.
  \item Introduce the intermediate vector $y=L^Tx$.
  \item Solve the lower-triangular system $Ly=b$ by forward substitution:
  \[
  y_i=\frac{b_i-\sum_{j=1}^{i-1}l_{ij}y_j}{l_{ii}},
  \qquad i=1,\ldots,n.
  \]
  \item Solve the upper-triangular system $L^Tx=y$ by back substitution:
  \[
  x_i=\frac{y_i-\sum_{j=i+1}^{n}l_{ji}x_j}{l_{ii}},
  \qquad i=n,n-1,\ldots,1.
  \]
  \item The resulting vector $x$ is the unique solution of $Ax=b$.
\end{enumerate}

\section{Symmetric Positive-Definite Matrices}

\subsection{Definition and meaning}
\begin{enumerate}
  \item A real square matrix $A$ is \emph{symmetric} if $A=A^T$.
  \item It is \emph{positive definite} if $x^TAx>0$ for every nonzero $x\in\mathbb{R}^n$.
  \item The scalar $x^TAx$ is the quadratic form associated with $A$.
  \item The word \emph{positive} is crucial: the quadratic form is strictly positive in every nonzero direction.
  \item Geometrically, $f(x)=x^TAx$ is an upward-curving bowl with its unique minimum at $x=0$.
  \item Every eigenvalue of a symmetric positive-definite matrix is strictly positive.
  \item Equivalently, $A=Q\Lambda Q^T$, where $Q$ is orthogonal and every diagonal entry of $\Lambda$ is positive.
  \item Positive definiteness is stronger than positive semidefiniteness, which requires only $x^TAx\geq 0$.
  \item For example, $\begin{bmatrix}2&1\\1&2\end{bmatrix}$ is positive definite because its eigenvalues are $1$ and $3$.
  \item A symmetric positive-definite matrix is commonly called an SPD matrix.
\end{enumerate}

\subsection{Why positive definiteness matters}
\begin{enumerate}
  \item Since all eigenvalues of an SPD matrix are positive, none is zero.
  \item Therefore, $A$ is invertible and $Ax=b$ has exactly one solution for every $b$.
  \item Positive definiteness guarantees the existence and uniqueness of the Cholesky factorization $A=LL^T$ with positive diagonal entries in $L$.
  \item Each diagonal entry of $L$ is obtained by taking a square root.
  \item Positive definiteness guarantees that the quantity under every such square root is positive.
  \item It also ensures that no diagonal entry $l_{ii}$ is zero.
  \item Consequently, the divisions in forward and back substitution are well defined.
  \item Cholesky uses roughly half the arithmetic and storage of a general LU factorization.
  \item In optimization, an SPD Hessian makes a quadratic objective strictly convex and gives it a unique minimizer.
  \item Thus positive definiteness supplies existence, uniqueness, and an efficient method for solving the system.
\end{enumerate}

\end{document}
